\section{A projected Hermitian--Einstein equation}
\label{sec:projected-correction}

Assume throughout this section that \(W\) is reduced. The factors
\(L,M\) may still coincide. Let \(\Pi_a\) be the spectral projection
onto the complete \(k\)-dimensional small space \(E_a\) from
Proposition~\ref{prop:weak-spectrum}, and put \(Q_a=I-\Pi_a\).
We decrease \(a_0\) so that the small eigenvalues are at most \(c/4\)
and the complementary spectrum is at least \(c>0\).
For an \(H_a\)-Hermitian trace-free section \(v\), define
\begin{equation}\label{eq:centered-curvature}
 \mathcal M_a(v)=e^{v/2}
       \bigl(i\Lambda F_{H_ae^v}-\kappa I\bigr)e^{-v/2}.
\end{equation}
The conjugation is the fibre isometry from \(H_ae^v\) to \(H_a\);
thus \(\mathcal M_a(v)\) is \(H_a\)-Hermitian and trace-free.
It is not a holomorphic change of the bundle.

\subsection{Scaled norms and the complementary inverse}

Fix \(0<\alpha<1\) and a sufficiently small \(\eta>0\).
On \(B(x,4\eta\rho_a(x))\), the weight \(\rho_a\) is comparable
to \(\rho_a(x)\). Rescale this ball by \(\rho_a(x)^{-1}\).
The rescaled base metrics are uniformly elliptic with bounded
coefficients in every fixed smooth norm. There are also unitary
frames with uniformly bounded scaled connection coefficients.
On \(r\le C\sqrt a\), use the smooth normalized holomorphic
frame, the dilation \(z=\sqrt a\,y\), and unitary normalization
by \(H_a^{-1/2}\). The matrix and its inverse have uniformly
bounded derivatives on each fixed enlarged \(y\)-ball.
On \(r\ge\sqrt a\), use \eqref{eq:graded-unitary-map} and
fixed unitary line frames. Their connection difference has
derivatives \(O(a r^{-3-j})\) of order \(j\), by the explicit
formulas in the proof of Proposition~\ref{prop:residual}.
Rescaling by \(r\) gives the claimed bounds. The fixed exterior
is uniformly smooth. These descriptions overlap; none uses
the singular splitting at a point of \(\Sigma\).

Let \(X_a\) be the supremum of the rescaled \(C^{2,\alpha}\)
norms in these frames, including the global supremum norm.
Let \(Y_a\) be the supremum of the rescaled \(C^{0,\alpha}\)
norms of \(\rho_a(x)^2f\) on the same balls. Uniformly equivalent
intrinsic norms are
\begin{align}
 \|v\|_{X_a}\asymp{}&
 \|v\|_\infty+\sup_X\bigl(\rho_a|D_av|+
                            \rho_a^2|D_a^2v|\bigr)
 +\sup_x\rho_a(x)^{2+\alpha}
       [D_a^2v]_{\alpha;B(x,\eta\rho_a(x))},
                                                \label{eq:scaled-x}\\
 \|f\|_{Y_a}\asymp{}&
 \sup_X\rho_a^2|f|
 +\sup_x\rho_a(x)^{2+\alpha}
       [f]_{\alpha;B(x,\eta\rho_a(x))}.
                                                \label{eq:scaled-y}
\end{align}
The Hölder differences can be taken after local parallel transport.
For each fixed \(a>0\), these are ordinary Hölder spaces with
equivalent norms. We do not differentiate a choice of charts
when differentiating a section with respect to \(a\).
Proposition~\ref{prop:residual} and four-dimensional integration give
\begin{equation}\label{eq:source-norms}
 \|\Phi_a\|_{Y_a}+\|\Phi_a\|_{4/3}\le Ca,\qquad
 \|\rho_a^{-2}\|_1+\|\rho_a^{-2}\|_{4/3}\le C.
\end{equation}
In particular \(Y_a\) controls \(L^1\) and \(L^{4/3}\) uniformly.

\begin{lemma}\label{lem:weighted-inverse}
The low projection and complementary inverse satisfy
\begin{align}
 \|\Pi_af\|_{X_a}&\le C\|f\|_1,\qquad
 \|\Pi_a\|_{X_a\to X_a}+\|\Pi_a\|_{Y_a\to Y_a}\le C,
                                                    \label{eq:low-projection}\\
 \|G_af\|_{X_a}&\le C\ell(a)\|f\|_{Y_a},\qquad
 G_a=(\mathcal L_a|_{Q_a})^{-1}Q_a.
                                                    \label{eq:log-inverse}
\end{align}
\end{lemma}

\begin{proof}
For the rough connection operator, scalar heat domination gives
\begin{equation}\label{eq:rough-heat-domination}
 |e^{-t\mathcal L_a}f|(x)
 \le e^{-t\Delta_0/2}|f|(x).
\end{equation}
One may obtain this by applying the maximum principle to a
regularized norm; there is no curvature potential in
\(\frac12D_a^*D_a\). The semigroup and kernel forms are also
given in \cite[Theorem~2.13]{Guneysu12} and
\cite[Theorem~8.8(b), Remark~8.9]{Guneysu11}, respectively.
Complexification followed by restriction gives the real-bundle
statement. On our fixed closed four-manifold the scalar kernel
obeys
\[
 h_t(x,y)\le Ct^{-2}\exp\{-d(x,y)^2/(Ct)\},
 \qquad 0<t\le1,
\]
by the compact specialization of \cite[Corollary~3.1]{LiYau86}.
The factor \(1/2\) in \eqref{eq:rough-heat-domination} changes
only constants.

If \(e\) is an \(L^2\)-normalized low eigensection, evaluation
at heat time one gives \(\|e\|_\infty\le C\).
Uniform local estimates for its scaled elliptic equation
give \(\|e\|_{X_a}\le C\). Expansion in an orthonormal
basis of the fixed-dimensional space \(E_a\) proves
\eqref{eq:low-projection}, using \eqref{eq:source-norms}
and the boundedness of \(\rho_a\).

We record the kernel bound needed for the weighted inverse.
For \(0<t\le1\), the kernel of \(e^{-t\mathcal L_a}Q_a\)
is the full kernel minus its \(k\) low terms, and hence has
norm at most
\[
 Ct^{-2}\exp\{-d(x,y)^2/(Ct)\}+C.
\]
For \(t\ge1\), place a fixed positive-time smoothing factor at
both ends of the complementary semigroup. Heat domination
bounds their \(L^2\)-to-fibre norms, while the middle factor
is bounded by \(Ce^{-ct}\) by the gap. Integration in time
therefore yields the fibre operator-norm bound
\begin{equation}\label{eq:green-kernel}
 |G_a(x,y)|\le C(1+d(x,y)^{-2}),\qquad x\ne y.
\end{equation}
This is a bound in the current metrics, independent of any
singular frame.

The four-dimensional convolution estimate
\begin{equation}\label{eq:log-convolution}
 \sup_x\int_X\frac{1+d(x,y)^{-2}}{\rho_a(y)^2}\,\dd V_y
 \le C\ell(a)
\end{equation}
now gives the supremum bound for \(G_af\).
Here is a direct verification near a defect \(p\).
If \(d(x,p)=O(\sqrt a)\), the contribution within distance
\(O(\sqrt a)\) of \(p\) is bounded by rescaling, and the
remaining radial integral is bounded by
\(\int_{\sqrt a}^R r/(a+r^2)\,\dd r\).
If \(R_x=d(x,p)\) is larger, the radius-\(R_x/2\)
neighborhood of \(x\) has \(\rho_a^2\gtrsim a+R_x^2\)
and bounded contribution. The radius-\(R_x/2\) ball about
\(p\) contributes at most
\[
 CR_x^{-2}\int_0^{R_x/2}\frac{r^3}{a+r^2}\,\dd r\le C.
\]
The intervening comparable annulus away from \(x\) also
contributes a constant. On \(r>2R_x\), \(d(x,y)\gtrsim r\)
and the same logarithmic integral applies. The fixed
exterior is uniformly integrable. Summing over the finitely
many points proves \eqref{eq:log-convolution}.
Finally, the uniform scaled Schauder estimate for
\(\mathcal L_aG_af=Q_af\), together with
\eqref{eq:low-projection}, upgrades the supremum estimate to
\eqref{eq:log-inverse}.
\end{proof}

\subsection{Solving the projected equation}

\begin{theorem}\label{thm:projected-correction}
There are fixed \(K,a_0>0\) such that for \(0<a\le a_0\)
there is a unique \(v_a\in Q_aX_a\) in the ball
\(\|v_a\|_{X_a}\le Ka\ell(a)\) satisfying
\begin{equation}\label{eq:projected-equation}
 Q_a\mathcal M_a(v_a)=0.
\end{equation}
The metric \(\widehat H_a=H_ae^{v_a}\) is smooth and has
the prescribed determinant. Moreover
\begin{align}
 \|v_a\|_{X_a}&\le Ca\ell(a),&
 \|D_av_a\|_2&\le Ca,&
 \|\Phi_{\widehat H_a}\|_\infty&\le Ca,
                                                   \label{eq:corrected-bounds}\\
 \widehat Z_a&\sim av_0,&
 \widehat{\mathcal E}_a&=2\pi^2\sigma a^2+o(a^2).
                                                   \label{eq:corrected-energy}
\end{align}
Uniqueness is asserted only in the indicated ball and slice.
\end{theorem}

\begin{proof}
The exact finite metric-change formula is
\[
 F_{H_ae^v}=F_{H_a}
   +\db_{\End V}\bigl(e^{-v}D'_a(e^v)\bigr).
\]
Differentiating \eqref{eq:centered-curvature} at zero gives
\[
 i\Lambda\db D'_av+\tfrac12[v,\Phi_a]
 =\tfrac12\bigl((D'_a)^*D'_a+(D''_a)^*D''_a\bigr)v
 =\mathcal L_av.
\]
This is the Kähler identity on endomorphism-valued zero-forms:
the centering commutator cancels half the difference between
the two Dolbeault Laplacians. See
\cite[(3.2.39)--(3.2.41), (6.3.16)--(6.3.17)]{Kobayashi87}
for the commutators and curvature variation.
In particular the real linearization is valid even though
\(\Phi_a\) is not central.

Write
\begin{equation}\label{eq:nonlinear-remainder}
 \mathcal M_a(v)=\Phi_a+\mathcal L_av+N_a(v).
\end{equation}
On a fixed sufficiently small \(X_a\) ball, the finite
change formula and the uniform scaled coefficients imply
\begin{equation}\label{eq:quadratic-remainder}
 \|N_a(v)-N_a(w)\|_{Y_a}
 \le C(\|v\|_{X_a}+\|w\|_{X_a})\|v-w\|_{X_a}.
\end{equation}
Indeed, in a unit rescaled chart the centered curvature is
a smooth matrix expression of order two with uniformly
bounded coefficients. After subtracting its value and
first derivative, every term contains at least two factors
from \(v\) and its derivatives, and at most two total spatial
derivatives. The \(C^{0,\alpha}\) product estimate proves
\eqref{eq:quadratic-remainder}; multiplication by \(\rho_a^2\)
converts physical curvature to that expression. Terms
involving \(vD_a^2v\) are included in this estimate.

On \(Q_aX_a\), equation~\eqref{eq:projected-equation} becomes
\[
 v=-G_a\bigl(\Phi_a+N_a(v)\bigr).
\]
By \eqref{eq:source-norms}, \eqref{eq:log-inverse} and
\eqref{eq:quadratic-remainder}, its right side has norm at most
\(C\ell(a)(a+\|v\|_{X_a}^2)\). Its Lipschitz constant
on the radius-\(Ka\ell(a)\) ball is at most \(CKa\ell(a)^2\).
Choose \(K\) larger than twice the linear bound and then
decrease \(a_0\). The contraction theorem gives the claimed
solution and local uniqueness. For fixed \(a\), the right
side of the unprojected equation is a linear combination
of smooth low eigensections. Its principal coefficients
are a small perturbation of those of \(\mathcal L_a\).
Quasilinear elliptic regularity makes \(v_a\) smooth.
Hermitian symmetry gives positivity of \(H_ae^{v_a}\),
and \(\tr v_a=0\) fixes its determinant.

The remaining centered curvature is
\[
 \mathcal M_a(v_a)=\Pi_a\Phi_a+\Pi_aN_a(v_a).
\]
The \(L^1\)-to-\(X_a\) projection estimate gives its supremum
norm at most \(C(a+a^2\ell(a)^2)\le Ca\). Centering preserves
the fibre norm, so this proves the curvature estimate in
\eqref{eq:corrected-bounds}.

For the stronger gradient bound, pair the equation with
\(v_a\). The \(L^{4/3}\) bound in \eqref{eq:source-norms},
the same bound for the quadratic remainder, the uniform
Sobolev inequality \eqref{eq:uniform-sobolev} and the
complementary gap give
\[
 \mathfrak q_a(v_a,v_a)
 \le\|\Phi_a+N_a(v_a)\|_{4/3}\|v_a\|_4
 \le Ca\,\mathfrak q_a(v_a,v_a)^{1/2}.
\]
Here \(\|N_a(v_a)\|_{4/3}\le Ca^2\ell(a)^2\), and on
\(Q_a\) the \(L^2\) term in the Sobolev estimate is controlled
by \(\mathfrak q_a\). Thus \(\|D_av_a\|_2\le Ca\).

The relative metric has supremum logarithm \(O(a\ell(a))\);
hence \(\widehat Z_a/Z_a=1+O(a\ell(a))\). For the energy,
the finite connection-change formula gives pointwise
\[
 |(D'_{\widehat H_a}-D'_{H_a})\iota|
 \le C\sqrt{u_a}\,|D_av_a|.
\]
Since \(u_a\le a\), its \(L^2\) norm is \(O(a^{3/2})\).
Also \(\|D'_{H_a\otimes h_L^{-1}}\iota\|_2=O(a)\).
Comparing the squared norms and the two fibre metrics gives
\[
 |\widehat{\mathcal E}_a-\mathcal E_a|
 \le Ca^{5/2}+Ca^3+Ca^3\ell(a)=o(a^2).
\]
The balanced identities \eqref{eq:balanced-static} now prove
\eqref{eq:corrected-energy}.
\end{proof}

The correction uses the whole small spectral space, including all
three directions when the line factors repeat. The remaining
curvature is not yet identified with the derivative of the metric
family. That identification requires estimates with respect to
the parameter, which we prove next.
